\chapter{Hasil dan Pembahasan}
\begin{tcolorbox}[title=Daftar Anggota Kelompok]
    \begin{enumerate}
        \itemsep0em
        \item Mhs 1, NIM: 
        \item Mhs 2, NIM: 
        \item Mhs 3, NIM:
    \end{enumerate}
\end{tcolorbox}
\section{Entity Relationship Diagram}
Dari analisa atas proses bisnis dan batasan yang diberikan, pada bagian ini dapat dijelaskan beberapa hal, diantaranya: \textit{mapping cardinality}, atribut dan kunci yang relevan, dsb.

Diagram ER untuk Basis Data Akademik ditunjukkan pada Gambar \ref{gb:1}
\begin{figure}[!h]
    \centering
    \includegraphics[width=.45\linewidth]{logougm.png}
    \caption{Skema ERD Basis Data Akademik}
    \label{gb:1}
\end{figure}
\section{Skema Tabel-Tabel Dalam Basis Data}
\begin{tcolorbox}[title=Skema Tabel dalam Basis Data Akademik]
\begin{itemize}
    \itemsep0em
     \item \texttt{mahasiswa(\dots)}
     \item \texttt{dosen(\dots)}
     \item dst.
\end{itemize}  
\end{tcolorbox}
\section{Pertanyaan Bisnis Dan Penyelesaian Dengan SQL}
\subsection{Soal nomor 1}
\begin{tcolorbox}
  \textit{Daftar nama mata kuliah pilihan dan mata kuliah yang memprasyaratinya (cukup info nama dan kode matakuliah saja).  Matakuliah yang tidak memiliki prasyarat tetap ditampilkan!}  
\end{tcolorbox}
\begin{enumerate}
    \itemsep0em
    \item Tabel yang diperlukan untuk menjawab query:
    \item Kiriteria join (bila ada)
    \item Perintah SQL sesuai pertanyaan
    \begin{lstlisting}
SELECT 
    FROM
    \end{lstlisting}
    \item Alternatif perintah SQL (bila ada) disertai penjelasan memadai
\end{enumerate}

\subsection{Soal nomor 2}
\begin{tcolorbox}   
Daftar Matakuliah pilihan yang diambil oleh mahasiswa Ilmu Komputer angkatan '2024' pada semester ganjil tahun akademik 2026/2027!
\end{tcolorbox}
\begin{enumerate}
    \itemsep0em
    \item Tabel yang diperlukan untuk menjawab query:
    \item Kiriteria join (bila ada)
    \item Perintah SQL sesuai pertanyaan
    \begin{lstlisting}
SELECT 
    FROM
    \end{lstlisting}
    \item Alternatif perintah SQL (bila ada) disertai penjelasan memadai
\end{enumerate}

\subsection{Soal nomor 3}
\begin{tcolorbox}
    Daftar Matakuliah pilihan yang favorit untuk mahasiswa Ilmu Komputer, yaitu tiga matakuliah dengan pengambil terbanyak dalam tiga tahun terakhir!
\end{tcolorbox}
\begin{enumerate}
    \itemsep0em
    \item Tabel yang diperlukan untuk menjawab query:
    \item Kiriteria join (bila ada)
    \item Perintah SQL sesuai pertanyaan
    \begin{lstlisting}
SELECT 
    FROM
    \end{lstlisting}
    \item Alternatif perintah SQL (bila ada) disertai penjelasan memadai
\end{enumerate}

\subsection{Soal nomor 4}
\begin{tcolorbox}
Daftar mahasiswa yang memiliki Indeks Prestasi Semester ganjil tahun 2025/2026 tidak kurang dari 3,25! (pahami perbedaan antara IP Semester dan IP Kumulatif)
\end{tcolorbox}
\begin{enumerate}
    \itemsep0em
    \item Tabel yang diperlukan untuk menjawab query:
    \item Kiriteria join (bila ada)
    \item Perintah SQL sesuai pertanyaan
    \begin{lstlisting}
SELECT 
    FROM
    \end{lstlisting}
    \item Alternatif perintah SQL (bila ada) disertai penjelasan memadai
\end{enumerate}

\subsection{Soal nomor 5}
\begin{tcolorbox}
    Daftar mahasiswa dan IP Kumulatif yang diperolehnya saat ini! (matakuliah yang pernah diulang cukup dihitung sekali) 
\end{tcolorbox}
\begin{enumerate}
    \itemsep0em
    \item Tabel yang diperlukan untuk menjawab query:
    \item Kiriteria join (bila ada)
    \item Perintah SQL sesuai pertanyaan
    \begin{lstlisting}
SELECT 
    FROM
    \end{lstlisting}
    \item Alternatif perintah SQL (bila ada) disertai penjelasan memadai
\end{enumerate}
